% =========================================================
% Asymptotic Notation: Little-o Language
% =========================================================

\subsection*{Little-o Notation}
\label{subsec:little-o-notation}

Asymptotic notation separates the main term of an expression from the part
that becomes negligible in a limiting process. The notation $f=o(g)$ means
that $f$ is smaller than every fixed multiple of $g$ near the limiting point.
It does not say that $f$ is zero. It says that $f$ is negligible relative to
the chosen scale $g$.

\begin{definition}[Little-o at a Point]
\label{def:little-o-at-a-point}
Let $a\in\mathbb{R}$, let $f$ and $g$ be real-valued functions defined on a
punctured neighbourhood of $a$, and suppose $g(x)\neq 0$ for all $x$ sufficiently
close to $a$ with $x\neq a$. We write
\[
  f(x)=o(g(x))\qquad (x\to a)
\]
if and only if for every $\varepsilon>0$ there exists $\delta>0$ such that
\[
  0<|x-a|<\delta
  \quad\Longrightarrow\quad
  |f(x)|\leq \varepsilon |g(x)|.
\]
\end{definition}

\begin{remark*}[Standard quantified statement]
\[
\begin{aligned}
f(x)=o(g(x))\ (x\to a)
\Longleftrightarrow
&\forall\varepsilon>0\;\exists\delta>0\;\forall x\\
&\bigl(0<|x-a|<\delta
  \Rightarrow |f(x)|\leq \varepsilon |g(x)|\bigr).
\end{aligned}
\]
\end{remark*}

\begin{remark*}[Interpretation]
The scale $g$ is the unit of comparison. The assertion $f=o(g)$ says that no
matter how small a tolerance multiple $\varepsilon g$ is chosen, $f$ eventually
fits inside that multiple.
\end{remark*}

\NoLocalDependencies

\begin{definition}[Little-o at Infinity]
\label{def:little-o-at-infinity}
Let $f$ and $g$ be real-valued functions defined on some interval
$(M,\infty)$, and suppose $g(x)\neq 0$ for all sufficiently large $x$. We write
\[
  f(x)=o(g(x))\qquad (x\to\infty)
\]
if and only if for every $\varepsilon>0$ there exists $R>0$ such that
\[
  x>R
  \quad\Longrightarrow\quad
  |f(x)|\leq \varepsilon |g(x)|.
\]
\end{definition}

\begin{remark*}[Standard quantified statement]
\[
\begin{aligned}
f(x)=o(g(x))\ (x\to\infty)
\Longleftrightarrow
&\forall\varepsilon>0\;\exists R>0\;\forall x\\
&\bigl(x>R\Rightarrow |f(x)|\leq \varepsilon |g(x)|\bigr).
\end{aligned}
\]
\end{remark*}

\begin{remark*}[Interpretation]
The point version and the infinity version have the same logical shape. Only
the neighbourhood changes: small punctured intervals around $a$ are replaced
by tails $(R,\infty)$.
\end{remark*}

\NoLocalDependencies

\begin{definition}[Increment Little-o]
\label{def:increment-little-o}
Let $r$ be a real-valued function defined for all sufficiently small nonzero
$h$. We write
\[
  r(h)=o(h)\qquad (h\to 0)
\]
if and only if for every $\varepsilon>0$ there exists $\delta>0$ such that
\[
  0<|h|<\delta
  \quad\Longrightarrow\quad
  |r(h)|\leq \varepsilon |h|.
\]
\end{definition}

\begin{remark*}[Standard quantified statement]
\[
r(h)=o(h)\ (h\to 0)
\Longleftrightarrow
\forall\varepsilon>0\;\exists\delta>0\;\forall h\,
\bigl(0<|h|<\delta\Rightarrow |r(h)|\leq \varepsilon |h|\bigr).
\]
\end{remark*}

\begin{remark*}[Interpretation]
This is the form used in first-order approximation. A remainder $r(h)$ is
smaller than the increment $h$ itself, so the quotient $r(h)/h$ tends to zero.
\end{remark*}

\begin{dependencies}
\begin{itemize}
  \item \hyperref[def:little-o-at-a-point]{Little-o at a Point}
\end{itemize}
\end{dependencies}

\begin{proposition}[Little-o Quotient Characterisation]
\label{prop:little-o-quotient-characterization}
Let $a\in\mathbb{R}$, and let $f$ and $g$ be real-valued functions defined on a
punctured neighbourhood of $a$, with $g(x)\neq 0$ for $x$ sufficiently close to
$a$ and $x\neq a$. Then
\[
  f(x)=o(g(x))\ (x\to a)
  \quad\Longleftrightarrow\quad
  \frac{f(x)}{g(x)}\to 0\ (x\to a).
\]

\smallskip\noindent
\hyperref[prf:little-o-quotient-characterization]{Go to proof.}
\end{proposition}

\begin{remark*}[Standard quantified statement]
\[
\left[\forall\varepsilon>0\;\exists\delta>0\;
0<|x-a|<\delta\Rightarrow |f(x)|\leq\varepsilon |g(x)|\right]
\Longleftrightarrow
\frac{f(x)}{g(x)}\to 0.
\]
\end{remark*}

\begin{remark*}[Interpretation]
The quotient form is often the fastest test for little-o. The epsilon form is
often the safest proof form because it avoids manipulating a quotient until the
nonvanishing of the scale has been checked.
\end{remark*}

\begin{dependencies}
\begin{itemize}
  \item \hyperref[def:little-o-at-a-point]{Little-o at a Point}
\end{itemize}
\end{dependencies}

\begin{proposition}[Sum Rule for Little-o]
\label{prop:little-o-sum-rule}
If $f_1(x)=o(g(x))$ and $f_2(x)=o(g(x))$ as $x\to a$, then
\[
  f_1(x)+f_2(x)=o(g(x))\qquad (x\to a).
\]

\smallskip\noindent
\hyperref[prf:little-o-sum-rule]{Go to proof.}
\end{proposition}

\begin{remark*}[Standard quantified statement]
\[
f_1=o(g)\land f_2=o(g)
\Longrightarrow
f_1+f_2=o(g).
\]
\end{remark*}

\begin{remark*}[Interpretation]
Negligible errors can be added. The proof is the usual epsilon-budget split:
make each error smaller than half the final tolerance.
\end{remark*}

\begin{dependencies}
\begin{itemize}
  \item \hyperref[def:little-o-at-a-point]{Little-o at a Point}
\end{itemize}
\end{dependencies}

\begin{proposition}[Bounded Factor Rule for Little-o]
\label{prop:little-o-bounded-factor-rule}
If $f(x)=o(g(x))$ as $x\to a$ and $m$ is bounded near $a$, then
\[
  m(x)f(x)=o(g(x))\qquad (x\to a).
\]

\smallskip\noindent
\hyperref[prf:little-o-bounded-factor-rule]{Go to proof.}
\end{proposition}

\begin{remark*}[Standard quantified statement]
\[
f=o(g)\land \exists M>0\;\exists\eta>0\;\forall x\,
\bigl(0<|x-a|<\eta\Rightarrow |m(x)|\leq M\bigr)
\Longrightarrow mf=o(g).
\]
\end{remark*}

\begin{remark*}[Interpretation]
A bounded multiplier cannot turn a negligible error into a main-order term. It
only changes the constant in the epsilon estimate.
\end{remark*}

\begin{dependencies}
\begin{itemize}
  \item \hyperref[def:little-o-at-a-point]{Little-o at a Point}
\end{itemize}
\end{dependencies}
